\underline{Lemma} 1.10. \underline{Let $(V,Q)$ be the quadratic module
which is the sum of two quadratic modules $(V_1,Q_1)$ and $(V_2,Q_2)$,
where $V_i$ is locally free of rank $2m_i$ and $Q_i$ is nondegenerate.
Let $V_i=W_i\oplus W'_i$ be a decomposition of $V_i$ into two totally
singular subspaces $(i=1,2)$.  Then the diagram}

\begin{center}
\begin{tikzcd}[column sep=huge,row sep=large]
C(Q) \arrow[rr,dash,"\sim"] \arrow[d,dash,"\wr"]
  & {} & C(Q_1)\otimes' C(Q_2) \arrow[d,equal] \\
\operatorname{End}(\bigwedge(W_1\oplus W_2))
  \arrow[r,phantom,"\simeq"]
  & \operatorname{End}(\bigwedge W_1\otimes\bigwedge W_2)
  \arrow[r,phantom,"\simeq"]
  & \operatorname{End}(\bigwedge W_1)\otimes'
    \operatorname{End}(\bigwedge W_2)
\end{tikzcd}
\end{center}

\underline{is commutative.}

%%%%%%%%%% source scan idx 76 | folio 69 %%%%%%%%%%

The verification is left to the reader.  One deduces at once from this
lemma that

\begin{equation}\tag{1.10.1}
e(W_1\oplus W_2)
=e(W_1)e(W_2)+(1-e(W_1))(1-e(W_2)).
\end{equation}

1.11. The following construction is borrowed from Micali and Villamayor.
If $p:Z\longrightarrow S$ is a double étale covering of $S$, then $Z$ may
be regarded as a $\mathbb Z/2$-torsor over $S$.  Denote the addition of
torsors by $*$; for a section $s$ of $Z$, let $e(s)$ be the idempotent of
$p_*\mathcal O_Z$ that equals one on $s$ and zero elsewhere.  With the
notation of 1.10, there is a unique isomorphism

\begin{equation}\tag{1.11.1}
\Spec(Z(Q_1))*\Spec(Z(Q_2))\longrightarrow\Spec(Z(Q)),
\end{equation}

whose underlying addition morphism is

\begin{equation*}
+:\Spec(Z(Q_1))\times_S\Spec(Z(Q_2))
\longrightarrow\Spec(Z(Q)),
\end{equation*}

and which, after every base change, satisfies

\begin{equation*}
e(s+t)=e(s)e(t)+(1-e(s))(1-e(t)).
\end{equation*}

If $s(W)$ is the section of $\Spec(Z(Q))$ such that $e(s(W))=e(W)$, then
formula (1.10.1) becomes

\begin{equation*}
s(W_1\oplus W_2)=s(W_1)+s(W_2).
\end{equation*}

\underline{Proposition} 1.12. \underline{Let $A$ be a field, let $V$ be
a vector space of dimension $2m>0$ over $A$, let $Q$ be a nondegenerate
quadratic form on $V$, and let $W_1,W_2$ be two totally isotropic
subspaces of $V$ of dimension $m$.  Then $e(W_1)=e(W_2)$ if and only if
$\dim(W_1/(W_1\cap W_2))$ is even.}

\footnote{The printed proposition says that $V$ has dimension $2m$ ``over
$V$''; the base field introduced in the same sentence is $A$.}

%%%%%%%%%% source scan idx 77 | folio 70 %%%%%%%%%%

(Cf. Bourbaki, \emph{Algèbre}, Ch.~9, \S6, Exercise~28.)  Let
$e_1,\ldots,e_m,f_1,\ldots,f_m$ be a basis of $V$ such that

\begin{enumerate}[label=\alph*),leftmargin=2.5em]
\item $Q(\sum_i x_i e_i+\sum_i y_i f_i)=\sum_i x_i y_i$;
\item $W_1$ is generated by the $e_i$;
\item $W_2$ is generated by $(e_i)_{1\leq i\leq s}$ and
$(f_i)_{s<i\leq m}$.
\end{enumerate}

Then $V$ is the sum of the $V_i=Ae_i+Af_i$, and

\begin{equation*}
W_\alpha=\bigoplus_i(W_\alpha\cap V_i),
\qquad \alpha=1,2.
\end{equation*}

The addition formula 1.11 therefore reduces the proof to the case
$\dim V=2$.  If $W_1=W_2$, the assertion is clear.  Otherwise, if
$W_1=Ae$ and $W_2=Af$ with $\Phi(e,f)=1$, then
$e(W_1)=fe$ and $e(W_2)=ef=1-fe$.

\medskip
2. \underline{Quadrics}

2.1. Let $k$ be an algebraically closed field and $V$ a vector space of
dimension $n+2$ over $k$.  A smooth quadric of dimension $n$ over $k$ is
a $k$-scheme $X$ isomorphic to the subscheme of $\mathbb P(V^*)$ defined
by the equation $Q=0$, for some ordinary quadratic form $Q$ over $k$.

\underline{Remark} 2.2. For $n=0,1,2$, $X$ is a smooth quadric of
dimension $n$ over $k$ if and only if

\begin{equation*}
\begin{array}{rcl}
n=0 &:& X\simeq\Spec(k)\amalg\Spec(k),\\[1ex]
n=1 &:& X\simeq\mathbb P^1_k,\\[1ex]
n=2 &:& X\simeq\mathbb P^1_k\times\mathbb P^1_k.
\end{array}
\end{equation*}
